\documentclass[11pt,a4paper]{article}
\usepackage[UTF8]{ctex}
\usepackage{geometry}
\geometry{left=2.5cm,right=2.5cm,top=2.5cm,bottom=2.5cm}
\usepackage{amsmath,amssymb,amsthm}
\usepackage{hyperref}
\usepackage{booktabs}
\usepackage{enumitem}
\usepackage{xltabular}
\usepackage{algorithm}
\usepackage{algpseudocode}
\hypersetup{colorlinks=true,linkcolor=blue,urlcolor=blue}

\newtheorem{definition}{定义}[section]
\newtheorem{theorem}{定理}[section]
\newtheorem{proposition}{命题}[section]
\newtheorem{remark}{注记}[section]

\title{基于生成论的网络安全分析能力调度系统\\
第1卷：防御方理论}
\author{李龙胜}
\date{\today}

\begin{document}

\maketitle

\begin{center}
\textit{
这是给防御者的手册。\\
1.1 告诉你最优分析率怎么算。\\
1.2 教你零数学基础落地。\\
1.3 为你展示整套理论的数学骨架。\\
选择你需要的深度，跳过你不需要的严谨。
}
\end{center}

\vspace{5mm}

\begin{abstract}
第1卷聚焦防御方如何用生成论的择优公式裁定告警的最优分析率。
1.1 建立告警分级裁定的数学模型，
定义漏报损失和分析成本，
推导最优分析率的解析条件，
给出完整的算法流程和伪代码。
1.2 提供面向一线工程师的实施指南，
包括台账建立、参数估算、动态反馈规则和快速启动方案。
1.3 将整个安全运营过程形式化为生成步系统的实例，
验证四条公设的满足性，并证明贪心最优性定理。
全部内容保持CC BY 4.0协议公开。
\end{abstract}

\tableofcontents

\section{1.1 问题与模型：告警洪流中的最优选择}

\subsection{问题描述}

现代安全运营中心每天产生海量告警。
一个中型企业每日告警量可达3万条以上，
而负责分析的一线分析师通常只有3到5人。
传统做法只有两种：
要么逐条分析，但人力完全无法承受；
要么按经验规则抽样，但随意舍弃可能导致真正的攻击信号被淹没。

生成论的最佳记录粒度理论为这个问题提供了第三种选择：
不是完整记录一切，也不是随意舍弃，
而是通过择优公式，自动裁定每类告警的“最优分析率”。
系统在漏报损失和分析成本之间寻找最低总成本的平衡点。

\subsection{系统抽象}

整个安全运营过程被抽象为一个生成步系统：
\begin{itemize}
    \item \textbf{状态空间}：当前安全态势，包括资产状态、已知漏洞、正在进行的攻击活动。
    \item \textbf{生成步}：对一条告警的分析操作，分为三类：
    \begin{enumerate}[label=(\arabic*)]
        \item 忽略——采样率 $r = 0$，不消耗人力，漏报风险最高。
        \item 自动分析——采样率 $r \in (0,1)$，由SOAR剧本处理，消耗少量人力。
        \item 人工深度分析——采样率 $r = 1$，逐条分析，消耗大量人力，漏报风险最低。
    \end{enumerate}
    \item \textbf{记录成本}：包括状态偏差成本（漏报损失）和评估成本（分析耗时）。
    \item \textbf{容量上限}：每日可投入的总工时预算 $C$。
\end{itemize}

\subsection{成本函数}

\textbf{状态偏差成本} $V(r)$：度量因分析率 $r$ 不足而漏报攻击所造成的期望损失。
当 $r = 1$（完整分析）时，$V(1) = 0$。
当 $r < 1$ 时，漏报概率为 $1-r$，假设损失与该概率成正比：
\begin{equation}\label{eq:V}
V(r) = \alpha \cdot A \cdot (1 - r),
\end{equation}
其中 $A$ 是该类告警中攻击成功时的最大损失，
$\alpha$ 是敏感度系数。

\textbf{评估成本} $K(r)$：度量以分析率 $r$ 处理告警所需的人力成本。
分析每一条告警需要平均工时 $t$，分析师时薪为 $s$，则：
\begin{equation}\label{eq:K}
K(r) = \kappa \cdot r,
\end{equation}
其中 $\kappa = s \cdot t$ 是单条告警的单位分析成本。

\textbf{总记录成本} $J(r)$：
\begin{equation}\label{eq:J}
J(r) = V(r) + K(r) = \alpha A (1 - r) + \kappa r.
\end{equation}

\subsection{最优分析率（线性模型）}

系统需要选择 $r \in [0,1]$ 使得总成本 $J(r)$ 最小。
$J(r)$ 是关于 $r$ 的线性函数：
\[
J(r) = \alpha A + (\kappa - \alpha A) r.
\]
因此：
\begin{itemize}
    \item 若 $\kappa < \alpha A$（分析成本低于期望损失），
          最优分析率 $r^* = 1$——完整分析。
    \item 若 $\kappa > \alpha A$（分析成本高于期望损失），
          最优分析率 $r^* = 0$——完全忽略。
    \item 若 $\kappa = \alpha A$，$r^*$ 在 $[0,1]$ 内任意。
\end{itemize}

\textbf{大白话解释}：
\begin{itemize}
    \item 如果可能损失大于分析成本（$\alpha A > \kappa$）：每条都要看（$r^*=1$）。
    \item 如果可能损失小于或等于分析成本（$\alpha A \le \kappa$）：完全忽略（$r^*=0$）。
\end{itemize}

\subsection{非线性扩展（进阶）}

真实安全运营中，攻击信号往往具有带宽特性——
攻击手法变化越快，需要的采样率越高。
引入信号采样的非线性模型更为精确。

设告警信号的有效带宽为 $B$，信号的最大动态范围为 $A$。
当分析率 $r$ 低于奈奎斯特频率时，产生混叠误差：
\begin{equation}\label{eq:V_nonlinear}
V(r) = \alpha A^2 \cdot \left( \frac{2B}{r} - 1 \right)^2, \quad r < 2B.
\end{equation}
评估成本仍为 $K(r) = \kappa r$。

最优分析率 $r^*$ 由方程
\begin{equation}\label{eq:optimal_r}
\frac{1-r}{r^4} = \frac{\kappa}{2\beta}
\end{equation}
在 $(0, 1]$ 内的唯一解决定，其中 $\beta = \alpha A^2 / (2B)^2$。
在实际SOC部署中，可先使用线性模型快速裁定，
随着历史数据的积累，逐步迁移至非线性模型。

\subsection{算法流程}

\textbf{预处理与特征提取}：对每条告警提取源IP、目标IP、端口、攻击类型、规则ID、资产重要性、历史误报率等。

\textbf{告警分类}：根据攻击类型和目标资产重要性，将告警归入 $K$ 个预定义类别。
例如：
\begin{itemize}
    \item 类别1：高危漏洞利用（针对核心服务器）。
    \item 类别2：中危漏洞扫描（针对普通服务器）。
    \item 类别3：低危信息收集（端口扫描）。
    \item 类别4：内部横向移动检测。
\end{itemize}

每类告警从台账中读取参数 $A, B, \alpha, \kappa$。

\textbf{最优分析率计算}：对每类告警 $i$：
\begin{enumerate}
    \item 计算 $\beta_i = \alpha_i A_i^2 / (2B_i)^2$（若使用非线性模型）。
    \item 求解方程 $\frac{1 - r_i}{r_i^4} = \frac{\kappa_i}{2\beta_i}$ 得到 $r_i^*$。
    \item 若使用线性模型，$r_i^* = 1$ 若 $\kappa_i < \alpha_i A_i$，否则 $r_i^* = 0$。
\end{enumerate}

\textbf{配额分配与总工时约束}：
设第 $i$ 类告警的预计日到达量为 $N_i$，每条平均分析时间为 $t_i$。
所需总工时：
\[
T = \sum_{i=1}^{K} r_i^* \cdot N_i \cdot t_i.
\]
若 $T \le C$，所有配额均被批准。
若 $T > C$，按 $\beta_i$（漏报损失大、分析收益高的类别）从高到低排列，
从最低 $\beta_i$ 的类别开始依次减少 $r_i^*$（每次减少 0.05），直至 $T \le C$。

\textbf{实时裁定}：对每条告警：
\begin{enumerate}
    \item 查找其类别 $i$ 及当前最优分析率 $r_i^*$。
    \item 生成随机数 $p \in [0,1]$。
    \item 若 $p \le r_i^*$：
    \begin{itemize}
        \item 若 $r_i^* \ge 0.8$，标记为人工深度分析。
        \item 否则，标记为自动分析（由 SOAR 剧本处理）。
    \end{itemize}
    \item 若 $p > r_i^*$：标记为忽略。
\end{enumerate}

\textbf{自适应反馈}：每周统计各类告警中真实攻击的发现率：
\begin{itemize}
    \item 若某类告警漏报导致安全事件，人工干预将 $A_i$ 乘以 2。
    \item 若某类告警连续四周零有效攻击，自动将 $A_i$ 乘以 0.9。
    \item 重新计算 $\beta_i$ 和 $r_i^*$。
\end{itemize}

\subsection{伪代码}

\begin{algorithm}[ht]
\caption{基于生成论的自动化告警分级裁定}
\begin{algorithmic}[1]
\Require 告警流，台账数据，分析师时薪，每日总工时预算 $C$
\Ensure 每条告警的分析动作（忽略/自动/人工）
\State 初始化：从数据库加载每类告警的 $A, B, \alpha, \kappa$
\Loop
    \State 预处理：解析告警，提取特征
    \State 分类：将告警归入预定义类别
    \For{每类告警 $i = 1$ to $K$}
        \State $\beta_i \gets \alpha_i \cdot A_i^2 / (2B_i)^2$
        \State 求解 $\frac{1 - r_i}{r_i^4} = \frac{\kappa_i}{2\beta_i}$ 得到 $r_i^*$（或使用线性模型）
    \EndFor
    \State 预测当日各类告警到达量 $N_i$
    \State $T \gets \sum_i r_i^* \cdot N_i \cdot t_i$
    \If{$T > C$}
        \State 按 $\beta_i$ 降序排列类别
        \For{类别 $i$ 从低 $\beta_i$ 到高 $\beta_i$}
            \While{$T > C$ 且 $r_i^* > 0$}
                \State $r_i^* \gets \max(0, r_i^* - 0.05)$
                \State 重新计算 $T$
            \EndWhile
        \EndFor
    \EndIf
    \For{每条告警}
        \State 取类别 $i$，$r \gets r_i^*$
        \State $p \gets \text{random}(0,1)$
        \If{$p \le r$}
            \If{$r \ge 0.8$}
                \State 标记：人工深度分析
            \Else
                \State 标记：自动分析
            \EndIf
        \Else
            \State 标记：忽略
        \EndIf
    \EndFor
    \State 记录分析结果和攻击发现情况
\EndLoop
\State 每周：评估漏报事件，若发生漏报则 $A_i \gets 2A_i$；若连续四周零有效攻击则 $A_i \gets 0.9A_i$
\State 每周：更新数据库中各类告警的 $A_i$，重新计算后续参数
\end{algorithmic}
\end{algorithm}

\subsection{诚实标注}

\begin{center}
\begin{xltabular}{\textwidth}{c|X|c}
\toprule
\textbf{命题} & \textbf{状态} & \textbf{层级} \\
\midrule
告警生成步系统抽象 &
抽象合理，跨SOC适用性待验证 &
II（框架已建立） \\
线性成本模型下的最优分析率 &
解析解已给出，边界条件明确 &
I（推导严格） \\
非线性成本模型下的最优分析率 &
方程已推导，依赖信号带宽假设 &
II（定性正确，参数需校准） \\
动态反馈调整 $A$ 的规则 &
启发式规则已建立，非从公设严格导出 &
III（启发式猜想，待最优化证明） \\
告警分类的完备性 &
依赖专家经验，可能存在未覆盖类型 &
II（框架完备，分类需持续更新） \\
\bottomrule
\end{xltabular}
\end{center}

\section{1.2 实施指南}

\subsection{核心思路：算一笔经济账}

我们不是追求“每条都看”，也不是“凭经验挑着看”。
而是用简单的数学逻辑，自动算出每类告警该花多少精力去查。

这个逻辑只有一句话：
\begin{quote}
\textbf{分析成本高于可能损失，就少查；可能损失高于分析成本，就多查。}
\end{quote}

\subsection{准备工作：建立成本台账}

新建一张 Excel 表（或数据库表），包含以下字段：

\begin{center}
\begin{tabular}{c|c|c|c}
\toprule
\textbf{字段} & \textbf{含义} & \textbf{示例值} & \textbf{如何获取} \\
\midrule
告警类别 & 攻击类型或规则名称 & “永恒之蓝漏洞利用” & 按SIEM规则分组 \\
$A$ & 漏掉一条告警可能损失（元） & 1000000 & 可参考攻防演练扣分表，后期自动修正 \\
$B$ & 每天平均告警数量（条） & 300 & 从日志统计 \\
$\alpha$ & 敏感度系数 & 1.0 & 初期设为1，根据情况调整 \\
$t$ & 分析一条告警耗时（小时） & 0.02（约1分钟） & 分析师记录平均处理时间 \\
$s$ & 分析师时薪（元/小时） & 100 & 公司薪酬数据 \\
\bottomrule
\end{tabular}
\end{center}

\subsection{参数快速估算}

\begin{itemize}
    \item \textbf{$A$（可能损失）}：初期建议分三档——
    \begin{itemize}
        \item 涉及核心服务器、数据库、域控的告警：$10^5 \sim 10^6$ 元。
        \item 涉及普通服务器、终端、业务系统的告警：$10^3 \sim 10^4$ 元。
        \item 已被WAF/防火墙自动阻断的告警：$10^1 \sim 10^2$ 元。
    \end{itemize}
    \item \textbf{$t$（分析耗时）}：让一线分析师记录一周内各类型告警的平均处理时间，取分钟数除以60。
    \item \textbf{$\alpha$（敏感度）}：初期统一设 1.0。后续如果发现某类告警特别容易被利用，可以手动调高。
\end{itemize}

\subsection{每日人力分配}

每天上午 8:00，脚本自动统计各类告警的预计到达量 $N_i$（取前7天均值）。
计算当天所需总工时：
\[
T = \sum_{i=1}^{K} r_i^* \cdot N_i \cdot t_i.
\]
若 $T$ 超出当日可投入总工时 $C$，则按以下优先级调整：
\begin{enumerate}
    \item 优先保证 $\alpha A$ 最大的告警类别得到完整的分析率。
    \item 降低 $\alpha A$ 最小的告警类别的分析率，每次减少 0.1，直到总工时不超过 $C$。
\end{enumerate}
调整后的分析率 $r_i^\dagger$ 即为当天该类别告警的实际抽样比例。

\subsection{动态调节：让台账越跑越准}

每周进行一次复盘，用实际数据修正初始估算。

\textbf{触发“翻倍”机制}：
如果某类告警漏报导致安全事件，则手动将台账中该类的 $A$ 值乘以 5。
下周起，该类的分析率会自动飙升。

\textbf{触发“衰退”机制}：
如果某类告警连续四周都生成过，但分析师复查后全部是误报，
则将该类的 $A$ 值乘以 0.9。
下周起，该类的分析率会自动下降。

\textbf{定期审计}：
每季度检查一次台账，确认 $A$ 和 $t$ 的值是否仍然合理。

\subsection{快速启动方案}

如果您不想等系统全部搭建完成，
可以先用以下方式测试算法：

\begin{enumerate}[label=(\arabic*)]
    \item 选取过去一个月内出现频率最高的 10 类告警。
    \item 在 Excel 中为这 10 类告警填好 $A$ 和 $t$ 的估计值（其他参数用默认值）。
    \item 用公式 $r^* = 1$ 若 $\alpha A > \kappa$，否则 $r^* = 0$ 手算出最优分析率。
    \item 比对过去一个月您的团队实际花了多少时间在这些告警上。
    \item 如果算法建议“忽略”的告警类型恰好是您也觉得“看了也白看”的，说明算法有效。
\end{enumerate}

\subsection{常见问题}

\textbf{这和设备的威胁评级有什么区别？}
设备评的是“这个攻击有多危险”，我们算的是“这个告警值不值得花时间去查”。
一个高危告警如果每次都自动被阻断，查它就是浪费人力。

\textbf{如果我漏掉了一条告警导致出事，谁负责？}
这套算法的逻辑是帮您把时间花在最可能出事的地方。
如果出事，出事的告警一定是被算法标注为“低分析率”的——
那您需要反思的是：为什么这个本该是高分析率的告警，您当初给它标了低损失值？

\textbf{需要什么样的技术栈？}
只需要三种东西：一张台账表（Excel或MySQL）、一个可以运行Python脚本的服务器、
以及安全设备的数据接口。不需要机器学习，不需要大数据平台。

\textbf{我的团队很小，这样做有意义吗？}
团队越小，越需要把力气用在刀刃上。小团队根本不可能处理所有告警。

\section{1.3 生成步系统实例化}

（本节将安全运营形式化为生成步系统，验证公设满足性。
如果对数学严格性不感兴趣，可以跳过本节。）

\subsection{状态空间}

安全运营状态由两部分组成：安全态势和剩余分析能力。

\begin{definition}[安全运营状态]
设 $\mathcal{S}_{\text{asset}}$ 为资产安全态势空间，
$\mathcal{R} \in [0, C]$ 为已消耗的工时累计量，$C$ 为每日总工时预算。
则安全运营状态为
\[
s = (\mathcal{S}_{\text{asset}}, \mathcal{R}) \in \mathcal{S}_{\text{SOC}}.
\]
\end{definition}

\subsection{生成步：告警分析操作}

告警分析操作 $U$ 属于三类之一：忽略、自动分析、人工深度分析。
在实际裁定中，这三类通过分析率 $r \in [0,1]$ 统一决定。

\subsection{记录成本函数}

记录成本由漏报损失和分析耗时组成：
\[
R(s, U) = 
\begin{cases}
\alpha_i A_i, & \text{若 } U = \text{忽略}, \\[4pt]
\kappa_i^{\text{auto}} + \alpha_i A_i (1 - \eta_{\text{auto}}), & \text{若 } U = \text{自动}, \\[4pt]
\kappa_i^{\text{manual}}, & \text{若 } U = \text{人工}.
\end{cases}
\]
当使用分析率 $r_i$ 对整类告警裁定时，
该类别期望记录成本为
\[
J_i(r_i) = (1 - r_i) \alpha_i A_i + r_i \bar{\kappa}_i.
\]

\subsection{容量约束}

每日累计消耗工时不得超过 $C$：
\[
\sum_{i=1}^{K} r_i \cdot N_i \cdot t_i \le C.
\]

\subsection{贪心择优与最优分析率}

在容量约束不饱和时，贪心择优给出的最优分析率为
\[
r_i^* = 
\begin{cases}
1, & \text{若 } \alpha_i A_i > \bar{\kappa}_i, \\[4pt]
0, & \text{否则}.
\end{cases}
\]
饱和时按 $\beta_i$ 降序削减。

\subsection{公理满足性}

本实例满足生成步系统的四条公理：
\begin{itemize}
    \item \textbf{过程性}：任何告警到达时，至少可执行忽略操作。
    \item \textbf{记录性}：任何分析操作均消耗不可逆的分析能力，最小记录量 $R_{\min} = \min(\kappa_i, \alpha_i A_i) > 0$。
    \item \textbf{够用就行原则}：每条告警的裁定遵循贪心择优。
    \item \textbf{观测自洽}：成本函数由台账参数确定，不同分析师使用相同台账时裁定结果一致。
\end{itemize}

\subsection{与贪心最优性定理的对接}

贪心最优性定理（第0卷）的条件在此实例中对应：
\begin{itemize}
    \item 单调进展量 $f(s)$ = 分析师工作时数累积，严格单调递增。
    \item 成本凸性：在忽略—自动—人工三级操作中，
          记录成本对分析率 $r$ 呈线性，而高阶的 $B$ 重标定成本
          引入了严格凸性，使贪心路径最优。
    \item 成本分解：$R = R_{\text{base}} + R_{\text{trans}}$，满足定理条件。
\end{itemize}

因此，安全运营告警分级裁定引擎是生成步系统的可靠实例，
贪心择优定理保证其最优性。

\subsection{诚实标注}

\begin{center}
\begin{xltabular}{\textwidth}{c|X|c}
\toprule
\textbf{命题} & \textbf{状态} & \textbf{层级} \\
\midrule
安全运营GS实例的公理满足性 &
已逐条验证 &
I \\
最优分析率贪心择优的最优性 &
在容量不饱和时严格等价于全局最优；
饱和时在凸性假设下保持最优 &
II \\
\bottomrule
\end{xltabular}
\end{center}

\vspace{5mm}
\noindent\textbf{文档信息}：
\begin{itemize}
    \item 作者：李龙胜
    \item 版本：v1.0（四卷体系重构第1卷）
    \item 许可：CC BY 4.0
    \item 前置依赖：第0卷
\end{itemize}

\end{document}